\documentclass[11pt]{article}
\usepackage{amsmath,amssymb,amsthm}
\usepackage[margin=1in]{geometry}
\newtheorem{theorem}{Theorem}
\newtheorem{proposition}[theorem]{Proposition}
\newtheorem{lemma}[theorem]{Lemma}
\newtheorem{corollary}[theorem]{Corollary}
\theoremstyle{remark}
\newtheorem{remark}[theorem]{Remark}
\newcommand{\Kac}{\operatorname{Kac}}

\title{Subleading coefficients of the fixed-charge recursion\\ at the second Jantzen layer}
\author{Notes for the colliding-poles paper (\texttt{paper.tex}, Section 7.2); draft, 2026-11-26}
\date{}

\begin{document}
\maketitle

\noindent
Notation of \texttt{paper.tex}: chains and chain functions
$c^{(k)}_{\lambda\to\mu,j}$ (Definition 7.9), the recursion at the fixed charge
(Theorem 7.14), the layer levels $\ell^\pm_k$, the Coulomb momenta $m^\pm_k$
and $d_x=(x^2-(p'-p)^2)/(4pp')$ (Section 8), and the top Jantzen layer
(Theorem 8.2).  Throughout, $b_\star^2=-p'/p$ with $p,p'\ge2$ coprime,
$\lambda=h_{r,s}(b_\star)$ is an interior weight with $n_1=p'r+ps<pp'$, and
$n_0=p'r-ps$.  For a positive label $g=(u,v)$ put
\[
 x_g=p'u-pv,\qquad y_g=p'u+pv,
\]
so that $h_g(b_\star)=d_{x_g}$ and $z_g(b_\star)=h_{u,-v}(b_\star)=d_{y_g}$:
$|x_g|$ is the momentum of the weight of $g$ and $|y_g|$ the momentum of the
weight it leads to.  The second-layer end points are
$\mu^\pm_2=\lambda+\ell^\pm_2$, with momenta $m^\pm_2=2pp'\pm n_0$.

\begin{lemma}[Derivatives of the Kac weights]\label{lem:hprime}
For every positive label $g$,
$h_g'(b_\star)=\frac{b_\star}{2p'^2}\bigl(p'^2-p^2-x_gy_g\bigr)$, where $'$ is
$d/db$.
\end{lemma}

\begin{proof}
From $h_{u,v}(b)=\frac14\bigl((b+b^{-1})^2-(ub+vb^{-1})^2\bigr)$,
$h'_{u,v}(b)=\frac12\bigl((1-u^2)b-(1-v^2)b^{-3}\bigr)$.  At $b_\star$,
$b_\star^{-3}=b_\star\,b_\star^{-4}=b_\star p^2/p'^2$, so
$h'_{u,v}(b_\star)=\frac{b_\star}{2p'^2}\bigl(p'^2(1-u^2)-p^2(1-v^2)\bigr)$, and
$p'^2u^2-p^2v^2=x_gy_g$.
\end{proof}

\begin{lemma}[Chains to the second layer]\label{lem:chains2}
The chains from $\lambda$ to $\mu^+_2$ are exactly
\[
 \gamma_1=\bigl((r,s),(p+r,p'-s)\bigr),\qquad
 \gamma_2=\bigl((p-r,p'-s),(r,2p'-s)\bigr),
\]
and the chains from $\lambda$ to $\mu^-_2$ are exactly
\[
 \gamma_1'=\bigl((r,s),(p-r,p'+s)\bigr),\qquad
 \gamma_2'=\bigl((p-r,p'-s),(2p-r,s)\bigr).
\]
The residues $R_g$ of all six labels are regular at $b=b_\star$, and the two
links of each chain satisfy $h'_{g_0}(b_\star)\neq h'_{g_1}(b_\star)$.
\end{lemma}

\begin{proof}
A label $g$ has weight $\nu$ exactly when $x_g=\pm x$ for the momentum $x$ of
$\nu$, and then it leads to the weight of momentum $|y_g|$, with level
$(y_g^2-x_g^2)/(4pp')$.  For $x=n_0$: since $\gcd(p,p')=1$, the solutions of
$p'u-pv=n_0$ in positive integers are $(r+ip,s+ip')$, $i\ge0$, and those of
$p'u-pv=-n_0$ are $((i+1)p-r,(i+1)p'-s)$, $i\ge0$; their levels are
$\ell^+_{2i+1}$ and $\ell^-_{2i+1}$.  By the table in Section 8, all of these
except $\ell^\pm_1$ exceed $\ell^\pm_2$, and $\ell^\pm_1<\ell^\pm_2$.  So a chain
to $\mu^\pm_2$ starts with $(r,s)$ (leading to $\mu^+_1$, momentum $n_1$) or
with $(p-r,p'-s)$ (leading to $\mu^-_1$, momentum $2pp'-n_1$); it is not a
single label, since $\ell^\pm_1\neq\ell^\pm_2$.

The labels of weight $\mu^+_1$ solve $p'u-pv=\pm n_1$: $(r+ip,ip'-s)$ and
$(ip-r,ip'+s)$ with $i\ge1$.  At $i=1$ they are $(p+r,p'-s)$, of level
$\ell^+_2-rs$, leading to momentum $2pp'+n_0$, and $(p-r,p'+s)$, of level
$\ell^-_2-rs$, leading to $2pp'-n_0$; for $i\ge2$ the level exceeds
$\ell^\pm_2-rs$.  Likewise the labels of weight $\mu^-_1$ solve
$p'u-pv=\pm(2pp'-n_1)$; the two of smallest level are $(r,2p'-s)$, of level
$\ell^+_2-(p-r)(p'-s)$, leading to $\mu^+_2$, and $(2p-r,s)$, of level
$\ell^-_2-(p-r)(p'-s)$, leading to $\mu^-_2$, and the others exceed both.  A
chain to $\mu^\sigma_2$ whose second label leads to $\mu^{-\sigma}_2$ would need
a third label of weight $\mu^{-\sigma}_2$ and level
$|\ell^+_2-\ell^-_2|=2|n_0|$.  Since $y_g-x_g=2pv$ and $y_g+x_g=2p'u$, one
has $y_g\equiv x_g\pmod{2p}$ and $y_g\equiv-x_g\pmod{2p'}$; as
$n_1\equiv n_0\pmod{2p}$ and $n_1\equiv-n_0\pmod{2p'}$, a label whose weight has
momentum $\equiv\pm n_0\pmod{2pp'}$ leads to a momentum $\equiv\pm
n_1\pmod{2pp'}$, and $|y_g|>|x_g|$.  For a label of weight $\mu^{-\sigma}_2$
(momentum $2pp'\pm n_0>n_1,2pp'-n_1$) this momentum is at least $2pp'+n_1$, so
the label leads to a weight of layer at least three and its level is at least
$\ell_3-\ell^{-\sigma}_2>\ell^\sigma_2-\ell^{-\sigma}_2=2|n_0|$.  This proves the
list.

A factor $ab+eb^{-1}$ of $\mathcal A_{u,v}$ vanishes at $b_\star$ exactly when
$(a,e)\in\mathbb Z(p,p')\setminus\{0\}$, which requires $u\ge p$ and $v\ge p'$
($a\in[1-u,u]$, $e\in[1-v,v]$).  None of the six labels has both, and
$\mathcal N_{u,v}(b,d)$ is a polynomial, so each $R_g$ is regular at
$b_\star$.  By Lemma~\ref{lem:hprime}, $h'_{g_0}-h'_{g_1}$ is a nonzero multiple
of $x_{g_1}y_{g_1}-x_{g_0}y_{g_0}$; the values are computed in the next proof
and are nonzero.
\end{proof}

\begin{proposition}[Second layer]\label{prop:second-layer}
For each sign,
\[
 c^{(1)}_{\lambda\to\mu^\pm_2,1}(b_\star)
 =\pm\frac{m^\pm_2}{n_0}\,c^{(2)}_{\lambda\to\mu^\pm_2,0}(b_\star),
 \qquad
 c^{(2)}_{\lambda\to\mu^\pm_2,0}(b_\star)
 =\mp\frac{n_0}{n_1}\,R_{r,s}(b_\star)\,R_{p\pm r,\,p'\mp s}(b_\star),
\]
and $c^{(2)}_{\lambda\to\mu^\pm_2,0}(b_\star)=\mathcal N^\pm_2/\gamma^\pm_2$.
Consequently, for $\ell^\pm_2\le N<\ell_3$ the simple-pole coefficient
$A^\star_{\lambda,1,N}$ of Theorem 7.14 contains the term
\[
 \pm\frac{m^\pm_2}{n_0}\,\frac{\mathcal N^\pm_2(d)}{\gamma^\pm_2}\,
 \partial_\Delta H_{N-\ell^\pm_2}(c_\star,\Delta,d)\big|_{\Delta=\mu^\pm_2},
\]
the top coefficient times the ratio of the momenta of $\mu^\pm_2$ and
$\lambda$.
\end{proposition}

\begin{proof}
Write $\delta_g(t)=h'_gt+O(t^2)$ with $h'_g=h'_g(b_\star)$ and put
$\rho_\gamma=R_{g_0}(b_\star)R_{g_1}(b_\star)$ for a chain
$\gamma=(g_0,g_1)$.  By Lemma~\ref{lem:chains2},
\[
 P^{(k)}_\gamma\frac{\delta_{g_1}^j}{j!}
 =\frac{\rho_\gamma\,h'^{\,k-1}_{g_0}h'^{\,j}_{g_1}}{j!\,(h'_{g_0}-h'_{g_1})}\,
  t^{k+j-2}+O(t^{k+j-1}).
\]
For $(k,j)=(1,0)$ the sum over the two chains is $c^{(1)}_{\lambda\to\mu,0}$,
which is regular at $t=0$ (Lemma 7.13); so
$\sum_\gamma\rho_\gamma/(h'_{g_0}-h'_{g_1})=0$.  For $k+j=2$ the value at $t=0$
is
$c^{(k)}_{\lambda\to\mu,j}(b_\star)
=\sum_\gamma\rho_\gamma h'^{\,k-1}_{g_0}h'^{\,j}_{g_1}/(j!\,(h'_{g_0}-h'_{g_1}))$.
By Lemma~\ref{lem:hprime}, with $K=p'^2-p^2$, $h'_g/(h'_{g_0}-h'_{g_1})
=(K-x_gy_g)/(x_{g_1}y_{g_1}-x_{g_0}y_{g_0})$.

\emph{Sign $+$, $m=m^+_2=2pp'+n_0$.}  The momenta are
$(x,y)=(n_0,n_1),(n_1,m)$ along $\gamma_1$ and
$(-n_0,2pp'-n_1),(n_1-2pp',m)$ along $\gamma_2$, so
$x_{g_1}y_{g_1}-x_{g_0}y_{g_0}$ equals $n_1(m-n_0)=2pp'n_1$ on $\gamma_1$ and
$(n_1-2pp')(m-n_0)=-2pp'(2pp'-n_1)$ on $\gamma_2$, both nonzero.  Regularity
gives $\rho_{\gamma_2}=\rho_{\gamma_1}(2pp'-n_1)/n_1$.  Then
\[
 c^{(2)}_{\lambda\to\mu^+_2,0}(b_\star)
 =\frac{\rho_{\gamma_1}}{2pp'n_1}\bigl[(K-n_0n_1)-(K+n_0(2pp'-n_1))\bigr]
 =-\frac{n_0}{n_1}\rho_{\gamma_1},
\]
\[
 c^{(1)}_{\lambda\to\mu^+_2,1}(b_\star)
 =\frac{\rho_{\gamma_1}}{2pp'n_1}\bigl[(K-n_1m)-(K-(n_1-2pp')m)\bigr]
 =-\frac{m}{n_1}\rho_{\gamma_1}.
\]

\emph{Sign $-$, $m=m^-_2=2pp'-n_0$.}  Along $\gamma'_1$ the momenta are
$(n_0,n_1),(-n_1,m)$, along $\gamma'_2$ they are
$(-n_0,2pp'-n_1),(2pp'-n_1,m)$; the differences are $-n_1(m+n_0)=-2pp'n_1$ and
$(2pp'-n_1)(m+n_0)=2pp'(2pp'-n_1)$.  Regularity gives
$\rho_{\gamma'_2}=\rho_{\gamma'_1}(2pp'-n_1)/n_1$, and the same computation gives
$c^{(2)}_{\lambda\to\mu^-_2,0}(b_\star)=(n_0/n_1)\rho_{\gamma'_1}$ and
$c^{(1)}_{\lambda\to\mu^-_2,1}(b_\star)=-(m/n_1)\rho_{\gamma'_1}$.

For the identification with Theorem 8.2 take $N$ with
$\ell^\pm_2\le N<\ell_3$ in~(7.8) at $k=2$.  The chain functions
$c^{(2)}_{\lambda\to\mu,j}$ with $\mu-\lambda<\ell_2$ end at $\mu^\pm_1$ and are
$R_g\delta_g^{1+j}/j!$, which vanish at $b_\star$; those ending at $\mu^\pm_2$
vanish at $b_\star$ for $j\ge1$ by the expansion above.  Since
$H_{N-\ell^\pm_2}(c_\star,\cdot)$ is regular at $\mu^\pm_2$
(Theorem 8.2), (7.8) reads
$A^\star_{\lambda,2,N}=\sum_\pm c^{(2)}_{\lambda\to\mu^\pm_2,0}(b_\star)
H_{N-\ell^\pm_2}(c_\star,\mu^\pm_2,d)$.  Comparing with~(8.1) at
$N=\min(\ell^+_2,\ell^-_2)$ (where $H_0=1$) and then at
$N=\max(\ell^+_2,\ell^-_2)$ gives $c^{(2)}_{\lambda\to\mu^\pm_2,0}(b_\star)
=\mathcal N^\pm_2/\gamma^\pm_2$.  The last assertion is~(7.8) at $k=1$.
\end{proof}

\begin{corollary}[Second-layer Jantzen norms]\label{cor:gamma2}
With $w^\pm_2$ normalized to $L_{-1}^{\ell^\pm_2}$-coefficient $1$,
\[
 \gamma^\pm_2=\mp\frac{n_1}{n_0}\,
 \frac{1}{\mathcal A_{r,s}(b_\star)\,\mathcal A_{p\pm r,\,p'\mp s}(b_\star)}.
\]
\end{corollary}

\begin{proof}
$R_g=\mathcal A_g\mathcal N_g$, and $\mathcal N^\pm_2=\mathcal N_{r,s}\mathcal
N_{p\pm r,p'\mp s}$ at $b_\star$ by the factorization after Theorem 8.2
(case $k=2j$, $j=1$).
\end{proof}

\begin{remark}
At $c=0$, $\lambda=0$ ($(p,p')=(2,3)$, $r=s=1$, $n_0=1$, $n_1=5$) the
corollary gives $\gamma^-_2=5/(\mathcal A_{1,1}\mathcal A_{1,4}(b_\star))
=5/(\frac12\cdot(-\frac3{1120}))=-\frac{11200}3$, the value found in the
example after Definition 7.9 of \texttt{paper.tex}.  The coefficients with
$k+j=2$ are the leading ones at the second layer; the remaining coefficient
$c^{(1)}_{\lambda\to\mu^\pm_2,0}(b_\star)$ depends on the first-order
expansions of $R_g$ and $\delta_g$ and is not a function of the momenta.
\end{remark}

\end{document}
